\documentclass[aspectratio=169,11pt]{beamer}
\usetheme{Madrid}

\usepackage[utf8]{inputenc}
\usepackage[T1]{fontenc}
\usepackage[magyar]{babel}
\usepackage{amsmath, amssymb}
\usepackage{tikz}
\usepackage{xcolor}

%%% Megjelenítéshez automatikusan kiegészített csomagok %%%
\usepackage{pgfplots}
\pgfplotsset{compat=1.18}
\usepackage{graphicx}
\usepackage{booktabs}
\usepackage{microtype}

\usepackage{hyperref}

\begin{document}

\title{Ismétlés nélküli permutáció}
\author{Varga Dávid}
\institute{Matematika}
\date{\today}

\begin{frame}
    \titlepage
\end{frame}

\begin{frame}{1. A lényeg}
    \begin{block}{Definíció}
        Képzeld el, hogy van pár teljesen különböző dolgod, és egyszerűen csak sorba akarod rendezni őket. Az \textbf{ismétlés nélküli permutáció} nem más, mint annak a kiszámítása, hogy hányféleképpen állíthatod sorba ezeket a dolgokat úgy, hogy mindegyiket pontosan egyszer használod fel.
    \end{block}
\end{frame}

\begin{frame}{2. Hogyan is működik?}
    Gondolj arra, hogy van három barát: Anna, Balázs és Csaba, és be akarnak ülni a moziba három egymás melletti székre. Hányféleképpen ülhetnek le?
    
    \begin{itemize}
        \item Az \textbf{első székre} még bárki leülhet (3 lehetőség).
        \item A \textbf{második székre} már csak a maradék két ember ülhet (2 lehetőség).
        \item A \textbf{harmadik székhez} érve már nincs választásunk (1 lehetőség).
    \end{itemize}

    \begin{alertblock}{Miért szorzunk?}
        Mert ezek a döntések egymásba ágyazódnak, mint egy fa ágai. A végeredmény: $3 \cdot 2 \cdot 1 = 6$ féle sorrend.
    \end{alertblock}
\end{frame}

\begin{frame}{3. Általános képlet és speciális esetek}
    \begin{block}{Definíció}
        $n$ különböző elem egy lehetséges sorrendjét az elemek egy \emph{ismétlés nélküli permutációjának} nevezzük. Az $n$ elem összes lehetséges sorrendjének a száma $P_n$.
    \end{block}

    $$P_n = n! = n \cdot (n - 1) \cdot (n - 2) \cdot \dots \cdot 1$$

    \begin{itemize}
        \item \textbf{Vigyázz, mikor használod!} Csak akkor működik, ha az elemek \textbf{különböznek}.
        \item \textbf{A nulla faktoriális:} $0! = 1$ (ha nulla dolgot kell sorbaraknod, azt 1-féleképpen teheted meg).
    \end{itemize}
\end{frame}

\begin{frame}{4. Dolgozat gyorsteszt}
    \begin{exampleblock}{3 pontos ellenőrzés}
        Amikor a dolgozatban meglátsz egy feladatot, futtasd le ezt a 3 pontos ellenőrzést:
        \begin{itemize}
            \item \textbf{Ugyanannyi a hely, mint az elem?} (Nem marad ki semmi?)
            \item \textbf{Számít a sorrend?} (A-B nem ugyanaz, mint B-A?)
            \item \textbf{Minden elem különbözik?} (Nincsenek egyforma elemek?)
        \end{itemize}
        Ha mindháromra \textbf{IGEN} a válasz, a képleted: $n!$.
    \end{exampleblock}
\end{frame}

\begin{frame}{5. Tipikus feladatok: Alapfeladat és "ragasztós" trükk}
    \begin{block}{1. Alapfeladat (Sorrend a moziban)}
        \textit{Egy 5 fős társaság moziba megy. Hányféleképpen ülhetnek le egymás mellé egy 5 székből álló sorban?}
        
        \textbf{Megoldás:} $5! = 5 \cdot 4 \cdot 3 \cdot 2 \cdot 1 = 120$.
    \end{block}

    \begin{block}{2. A "ragasztós" trükk (Összetartozó elemek)}
        \textit{5 fő leül 5 székre, de Béla és Cili egymás mellett akar ülni.}
        
        \textbf{Megoldás:} A párost összeragasztjuk (4 egység $\rightarrow$ $4!$), majd szorozzuk a belső helycserével ($2!$). Így: $4! \cdot 2! = 48$.
    \end{block}
\end{frame}

\begin{frame}{6. Kizárási és korlátozó feladatok}
    \begin{block}{3. Kizárási feladat}
        \textit{Hányféleképpen ülhet le 5 ember úgy, hogy Anna és Béla NEM ülhet egymás mellett?}
        
        \textbf{Megoldás:} Összes eset $-$ tiltott esetek $\rightarrow$ $120 - 48 = 72$.
    \end{block}

    \begin{block}{4. Fixált helyek (Számok képzése)}
        \textit{A 0, 1, 2, 3, 4 számjegyekből hány darab 5-jegyű szám képezhető?}
        
        \textbf{Megoldás:} Az 1. helyre nem jöhet 0 (4 lehetőség), a többi 4 helyre $4!$ féleképpen kerülnek a számok. $4 \cdot 4! = 96$.
    \end{block}
\end{frame}

\begin{frame}{7. Körmutáció, Rögzített pozíció, Paritás}
    \begin{block}{További típusok}
        \begin{itemize}
            \item \textbf{Ciklikus / Körmutáció:} $n$ ember kerek asztal körül: $(n-1)!$ (egy embert rögzítünk). 5 főre: $4! = 24$.
            \item \textbf{Kötött pozíció:} Leggyorsabb futó fixen az 1. helyen $\rightarrow$ $1 \cdot 5! = 120$.
            \item \textbf{Váltakozó elrendezés:} 3 fiú és 3 lány felváltva $\rightarrow$ $2 \cdot 3! \cdot 3! = 72$.
        \end{itemize}
    \end{block}
\end{frame}

\begin{frame}{Köszönöm a figyelmet!}
    \begin{center}
        \Large Kérdések?
    \end{center}
\end{frame}

\end{document}